%%%%%%%%%%%%%%%%%%%%%%%%%%%%%%%%%%%%%%%%%%%%%%%%%%%%%%%%%%%%%%%%%%%%%%%%

%%% LaTeX Template for AAMAS-2026 (based on sample-sigconf.tex)
%%% Prepared by the AAMAS-2026 Publication Chairs based on the version from AAMAS-2025. 

%%%%%%%%%%%%%%%%%%%%%%%%%%%%%%%%%%%%%%%%%%%%%%%%%%%%%%%%%%%%%%%%%%%%%%%%

%%% Start your document with the \documentclass command.

\documentclass[sigconf,anonymous]{aamas} 

%%% Load required packages here (note that many are included already).

\usepackage{balance} % for balancing columns on the final page
\usepackage{algorithm}
\usepackage{algorithmic}
\usepackage{booktabs}
\usepackage{multirow}
\usepackage{colortbl}
\usepackage{xcolor}
\usepackage[section]{placeins}
\usepackage{amsmath}
\DeclareMathOperator*{\argmax}{arg\,max}
\DeclareMathOperator*{\argmin}{arg\,min}
\usepackage{caption}
\usepackage{subcaption}
\usepackage{todonotes}
\usepackage{orcidlink}
\usepackage{hyperref}


%%%%%%%%%%%%%%%%%%%%%%%%%%%%%%%%%%%%%%%%%%%%%%%%%%%%%%%%%%%%%%%%%%%%%%%%

%%% AAMAS-2026 copyright block (do not change!)

\setcopyright{ifaamas}
\acmConference[AAMAS '26]{Proc.\@ of the 25th International Conference
on Autonomous Agents and Multiagent Systems (AAMAS 2026)}{May 25 -- 29, 2026}
{Paphos, Cyprus}{C.~Amato, L.~Dennis, V.~Mascardi, J.~Thangarajah (eds.)}
\copyrightyear{2026}
\acmYear{2026}
\acmDOI{}
\acmPrice{}
\acmISBN{}


%%%%%%%%%%%%%%%%%%%%%%%%%%%%%%%%%%%%%%%%%%%%%%%%%%%%%%%%%%%%%%%%%%%%%%%%

%%% == IMPORTANT ==
%%% Use this command to specify your submission number.
%%% In anonymous mode, it will be printed on the first page.

\acmSubmissionID{<<submission id>>}

%%% Use this command to specify the title of your paper.

\title[Counterfactual Gradient Alignment]{Counterfactual Gradient Alignment: Optimizing Directional Expert Supervision for Data-Efficient Learning}

%%% Provide names, affiliations, and email addresses for all authors.

\author{Arthur Pendragon}
\affiliation{
  \institution{Camelot Castle}
  \city{Camelot}
  \country{United Kingdom}}
\email{king.arthur@camelot.uk}

\author{Nimue}
\affiliation{
  \institution{The Lady's Lake}
  \city{Avalon}
  \country{United Kingdom}}
\email{lady.of.the.lake@avalon.uk}

%%% Use this environment to specify a short abstract for your paper.

\begin{abstract}
Supervised learning traditionally relies on passive data-label pairs, often failing to capture the rich geometric intuition that human experts possess about decision boundaries. This limitation is particularly acute in low-data regimes where models lack sufficient signal to generalize. We propose \textbf{Counterfactual Gradient Alignment (CGA)}, a framework that leverages expert-provided directions pointing towards counterfactual instances to supervise the model's gradient landscape. By penalizing positive directional derivatives along these vectors, we effectively teach the model the local topology of class transitions. We evaluate three novel loss formulations---ReLU, Softplus, and Sign-based variants---across a diverse suite of benchmarks including synthetic 2D datasets, image classification (MNIST), and sentiment analysis (IMDB, SST-2, SST-5). Our results demonstrate significant accuracy boosts in sparse data environments, such as a +7.17\% gain on IMDB with only 200 samples and a +14.5\% improvement on complex 2D circle boundaries. Crucially, we identify a profound synergy between CGA and active learning, showing that directional supervision is most impactful when applied to model-selected informative boundary samples, whereas it can act as noise on random data subsets.
\end{abstract}

%%% Use this command to specify a few keywords describing your work.
%%% Keywords should be separated by commas.

\keywords{Counterfactual Explanations, Gradient Supervision, Active Learning, Data Efficiency, Human-AI Interaction}

%%%%%%%%%%%%%%%%%%%%%%%%%%%%%%%%%%%%%%%%%%%%%%%%%%%%%%%%%%%%%%%%%%%%%%%%

%%% Include any author-defined commands here.         
\newcommand{\BibTeX}{\rm B\kern-.05em{\sc i\kern-.025em b}\kern-.08em\TeX}

%%%%%%%%%%%%%%%%%%%%%%%%%%%%%%%%%%%%%%%%%%%%%%%%%%%%%%%%%%%%%%%%%%%%%%%%

\begin{document}

%%% The following commands remove the headers in your paper. For final 
%%% papers, these will be inserted during the pagination process.

\pagestyle{fancy}
\fancyhead{}

%%% The next command prints the information defined in the preamble.

\maketitle 

%%%%%%%%%%%%%%%%%%%%%%%%%%%%%%%%%%%%%%%%%%%%%%%%%%%%%%%%%%%%%%%%%%%%%%%%

\section{Introduction}
Modern machine learning models, while highly capable, typically require vast amounts of data to learn reliable decision boundaries and generalise well to out-of-distribution examples \cite{hand_classifier, Challen231, souza2020challenges}. Humans are, by contrast, significantly more data-efficient. We utilize information beyond passive data-label pairs, applying structural priors and geometric intuition to generalize from a handful of examples. For instance, a human expert doesn't just recognize a specific digit as a `3'; they understand that changing a specific curve would turn it into an `8'. This directional intuition represents a latent supervision signal that is largely untapped in current supervised learning pipelines.

In this work, we propose \textbf{Counterfactual Gradient Alignment (CGA)}, a framework that converts this intuitive knowledge into a differentiable training signal. In CGA, an expert (be it a human or an automated oracle) provides a direction vector $\mathbf{d}$ indicating the shortest or most feasible path to a counterfactual instance---a point where the classification changes. We then constrain the model's gradients to align with this vector, ensuring that moving towards the counterfactual consistently reduces the model's confidence in the original class. This process effectively teaches the model the local geometry of the decision boundary, providing a much denser supervision signal than a simple categorical label.

We evaluate CGA across vision and NLP domains, covering both binary and multiclass tasks. We explore a family of loss functions---ReLU, Softplus, and Sign-based variants---and investigate the interaction between CGA and importance sampling (Active Learning). Our experiments demonstrate that CGA provides substantial gains in low-data regimes, particularly when paired with active selection of boundary samples. Specifically, we observe a +7.17\% gain on IMDB sentiment and a massive +14.5\% improvement on complex 2D boundaries.

A key discovery of our study is the profound synergy between CGA and active learning. We demonstrate that directional supervision is uniquely effective when applied to model-selected informative boundary samples. Conversely, applying directional alignment to random or "easy" samples can be counter-productive, highlighting that CGA is a high-precision regularizer that requires careful sample targeting. This suggests that CGA is best deployed in a "teaching" loop where the model proactively queries the expert for directional intuition on its most ambiguous boundaries.

\section{Related Work}
\label{related_work}

\subsection{Embedding human priors in model training}
One significant motivation of this work is to enable human annotators to provide domain knowledge to models to improve data efficiency. \citeauthor{rieger2020interpretations} \cite{rieger2020interpretations} demonstrated the ability to embed human priors using contextual decomposition explanation penalisation (CDEP). This method assumes that the explanation method is truthful and penalises model explanations where these do not align with an oracle. Similarly, Ross et al. embed human priors as explanations and develop a loss function which balances cross entropy with an additional function to penalise misaligned feature attributions \cite{ross2017right}.

\subsection{Training with counterfactuals}
Kaushik et al. \cite{Kaushik2020Learning} present a system for collecting counterfactually augmented data (CAD) by tasking human annotators to minimally edit individual observations to produce a counterfactual on a subset of the IMDB dataset. They show that incorporating counterfactuals within the training data as additional factual points improves generalisation. In our work, we take a different approach: we use counterfactuals not just as points, but as \emph{directional supervisors} for the model's derivative landscape.

\subsection{Learning from counterfactuals through gradient supervision}
\label{subsec:gradient_supervision}
Teney et al. \cite{teney2020learning} introduce "gradient supervision", which involves penalising the angle between the gradients of a neural network and the vector difference between counterfactual examples. Using a counterfactual, they define a loss function:
\begin{equation}
        \mathcal{L}_{GS} \left(\mathbf{g}_i,\hat{\mathbf{g}}_i \right)  = 
        1 - \langle\mathbf{g}_i,\hat{\mathbf{g}}_i \rangle /
        \left( \lVert \mathbf{g}_i \rVert \lVert \hat{\mathbf{g}}_i \rVert \right)
\end{equation}
where $\mathbf{g}_i$ denotes the model gradient and $\hat{\mathbf{g}}_i$ the counterfactual direction. While Teney et al. focus on angular alignment, our proposed loss functions focus on penalizing the directional derivative directly, exploring smoother landscapes (Softplus) and bounded signals (Sign) which we find to be more robust across diverse dataset topologies.

\subsection{Importance Sampling}
Importance sampling techniques in active learning prioritize samples that contribute most to model learning. Settles and Craven~\cite{settles2008active} demonstrated that variable annotation costs can significantly improve learning curves. He et al.~\cite{he2014active} combined uncertainty sampling with representativeness. Our work investigates the interaction between these sampling strategies and gradient-based supervision, showing that CGA depends uniquely on sample informativeness.

\section{Gradient-Based Learning}
\label{sec:gradient_learning}
In traditional supervised learning, we assume a dataset of random observations {\(\mathbf{x_i}\), \(y_i\)}. In CGA, we consider instead triplets {\(\mathbf{x_i}\), \(y_i\), $\mathbf{d_i}$}, where $\mathbf{d_i}$ defines a unit vector from a point $\mathbf{x_i}$ pointing towards a proximal example of the opposing class; a \emph{counterfactual}.

Formally, a counterfactual $\hat{\mathbf{x}}_i$ is the closest point to $\mathbf{x_i}$ for which the class changes. We propose that the direction towards this counterfactual should indicate a decrease in prediction probability for the true class $y$. We define the counterfactual direction as:
\begin{equation}
    \mathbf{d_i}=\frac{\hat{\mathbf{x}}_i - \mathbf{x_i}}{|\!|\hat{\mathbf{x}}_i - \mathbf{x_i}|\!|}. \label{eq:dir_unit}
\end{equation}
To measure for decreasing probability, we calculate the \textbf{Directional Derivative} $d$:
\begin{equation}
    d = \nabla_x f_{y}(x) \cdot \mathbf{d}_{i}
\end{equation}
where $f_y(x)$ is the model's score for the true class. We aim to minimize $d$ (make it negative). We evaluated three directional loss variants $\mathcal{L}_{d}$:

\begin{enumerate}
    \item \textbf{ReLU Loss}: $\mathcal{L}_{\text{ReLU}} = \max(0, d)$. Penalizes only positive directional derivatives.
    \item \textbf{Softplus Loss}: $\mathcal{L}_{\text{Softplus}} = \log(1 + e^{\beta d})$. A smooth variant that encourages a steep confidence drop-off even when the model is already "correct."
    \item \textbf{Sign Loss}: $\mathcal{L}_{\text{Sign}} = \tanh(\beta d) + 1$. Focuses purely on the sign of the derivative, providing a bounded signal robust to outliers.
\end{enumerate}

We combine the directional loss with standard cross entropy loss $\mathcal{L}_{\textrm{CE}}$ using a mixing parameter $\alpha \in [0,1]$:
\begin{align}\label{eq:lambda}
    \mathcal{L}_{\text{Total}} = (1-\alpha)\mathcal{L}_{\textrm{CE}} + \alpha \mathcal{L}_{d}
\end{align}

\section{Experimental Setup}
\label{sec:experimental_setup}

\subsection{Tasks and Datasets}
\paragraph{2D Synthetic:} XOR, Concentric Circles, and Two Moons. These tasks allow for clear visualization of decision boundary shifts under directional supervision.
\paragraph{Image Classification:} MNIST digits with three path types for counterfactuals: Raw pixel paths, PCA-transformed paths, and feasible paths generated by the FACE algorithm \cite{poyiadzi2020face}.
\paragraph{Text Sentiment:} IMDB (binary), SST-2 (binary), and SST-5 (5-class). These represent high-dimensional, discrete input spaces.

\subsection{Implementation and Training}
All models were implemented in \textbf{JAX/Flax}. For image classification, we used a CNN (\texttt{OptimizedMNISTConv}). For text, we used Bag-of-Words (BoW) neural networks with trainable embedding layers. Optimization was performed using the \textbf{AdamW} optimizer with a peak learning rate of $10^{-3}$ and weight decay of $10^{-4}$. We simulated data scarcity by subsetting training data to \{10, 20, 50, 100, 200\} samples.

\subsection{Active Learning Protocol}
To evaluate the impact of sample quality, we compared two data acquisition strategies: (1) \textbf{AL ON}: Entropy-based importance sampling with a diversity parameter to select boundary samples; (2) \textbf{AL OFF}: Purely random sample selection.

\section{Results}
\label{sec:results}

\subsection{NLP Benchmarks}
As shown in Table \ref{tab:nlp_results}, CGA achieves substantial gains on IMDB and SST-5 sentiment tasks. The Softplus variant was the most consistent performer, providing a +7.17\% gain on IMDB.

\begin{table}[h]
\centering
\caption{Validation Accuracy on IMDB and SST-5 sentiment benchmarks (AL ON, 5 Epochs).}
\label{tab:nlp_results}
\begin{tabular}{lcccc}
\toprule
\textbf{Dataset} & \textbf{Size} & \textbf{Baseline (CE)} & \textbf{Best CGA} & \textbf{Gain} \\ \midrule
IMDB & 200 & 55.94\% & \textbf{63.11\% (SP)} & \textbf{+7.17\%}
 \\IMDB & 100 & 51.43\% & \textbf{56.56\% (Sign)} & \textbf{+5.13\%}
 \\SST-5 & 100 & 53.06\% & \textbf{56.46\% (ReLU)} & \textbf{+3.40\%}
 \\bottomrule
\end{tabular}
\end{table}

\subsection{Synergy with Active Learning}
The most dramatic finding was the absolute dependency of CGA on informative samples. Without active selection, CGA performed worse than the baseline on all text tasks (Table \ref{tab:al_synergy}).

\begin{table}[h]
\centering
\caption{Performance Gain over Baseline at Size 100. CGA utility flips from positive to negative without Active Learning.}
\label{tab:al_synergy}
\begin{tabular}{lcc}
\toprule
\textbf{Dataset} & \textbf{AL ON (Boundary)} & \textbf{AL OFF (Random)} \\ \midrule
Gain (IMDB) & \textbf{+5.13\%} & -2.26\%
 \\Gain (SST-2) & \textbf{+1.36\%} & -2.72\%
 \\Gain (SST-5) & \textbf{+3.40\%} & -4.77\% \\bottomrule
\end{tabular}
\end{table}

\subsection{Complex Topologies: 2D Circles}
CGA's ability to regularize complex topologies was most evident in the Concentric Circles task, where directional information guided the model to form the correct circular boundary even when factual samples were sparse.

\begin{table}[h]
\centering
\caption{Validation Accuracy on Concentric Circles (Size 100).}
\label{tab:2d_circles}
\begin{tabular}{lcccc}
\toprule
\textbf{Dataset} & \textbf{Baseline (CE)} & \textbf{Best CGA} & \textbf{Variant} & \textbf{Gain} \\ \midrule
Circles & 57.50\% & \textbf{72.00\%} & Sign & \textbf{+14.50\%}
 \\bottomrule
\end{tabular}
\end{table}

\section{Discussion}
\label{sec:discussion}
The consistent superiority of the \textbf{Softplus} variant suggests that smooth, persistent gradient alignment is more effective than hard hinge penalties. By providing a non-zero loss even when the directional derivative is negative, Softplus ensures the model maintains a high-margin confidence drop toward the counterfactual boundary, leading to more robust decision regions.

The dependency on active learning reveals that directional supervision is a high-precision signal. For central, "easy" samples, the counterfactual direction may align with irrelevant features or conflict with the classification signal. However, for boundary samples where the model is uncertain, the directional signal provides critical structural information that categorical labels lack.

\section{Conclusions}
\label{sec:conclusions}
Counterfactual Gradient Alignment (CGA) allows models to inherit geometric intuition from experts, significantly improving generalization in data-scarce environments. Our work demonstrates that smooth directional losses combined with active boundary selection can achieve accuracy gains of over 14\%. Future research will investigate the application of CGA to large-scale generative models and multi-modal grounding.

\balance
\bibliographystyle{ACM-Reference-Format}
\bibliography{sample}

\appendix
\section{Counterfactual Path Generation}
On MNIST, we found that feasible paths generated by the FACE algorithm \cite{poyiadzi2020face} provide more useful directions than raw pixel-space vectors, especially as the training size increases (Table \ref{tab:mnist_paths}).

\begin{table}[h]
\centering
\caption{MNIST Accuracy Comparison across Path Types (Size 200).}
\label{tab:mnist_paths}
\begin{tabular}{lcccc}
\toprule
\textbf{Path Type} & \textbf{Baseline} & \textbf{Best Accuracy} & \textbf{Gain} \\ \midrule
Raw & 83.50\% & 84.52\% & +1.02\%
 \\PCA & 83.50\% & 84.56\% & +1.06\%
 \\FACE & 83.50\% & \textbf{84.78\%} & \textbf{+1.28\%}
 \\bottomrule
\end{tabular}
\end{table}

\end{document}
%%%%%%%%%%%%%%%%%%%%%%%%%%%%%%%%%%%%%%%%%%%%%%%%%%%%%%%%%%%%%%%%%%%%%%%%
